\section{Conclusion}
\label{sec:conclusion}

Our study demonstrates that redesigning the final layers of an Advanced LIGO-style mirror coating can improve its tolerance to random thickness errors. When we redesigned the top four layers using perfect measurements, the median tolerable error across 32 designs increased from 1.14\% to 1.71\%, about a 1.4-fold improvement. The rescue improved 30 of the 32 designs. The change in thermal noise was negligible. Tests with six and eight redesigned layers showed that there is room for further improvement.

However, the method depends strongly on measurement accuracy. It worked best when the deposited layers were measured more accurately than they were made. When measurement error became as large as the deposition error, the rescue no longer helped, and when it was larger, the rescue did harm. Using it only on mirrors predicted to fail could not solve this problem completely.

The method was also much less effective against shared errors, where all layers become too thick or too thin by a similar amount. At 1\% shared error, even after redesigning up to 12 layers, only two of the 32 designs reached a 95\% pass rate. Such errors may require correcting the deposition process itself rather than redesigning only the final layers.

Overall, the rescue offers a way to reduce the effect of random layer-to-layer thickness errors while keeping thermal noise nearly unchanged. Its practical value will depend on the accuracy of the coating measurements and the ability to identify the type of error affecting the coating. Further simulations and experimental tests are needed before the method can be applied to real mirror coatings.
